\pdfgentounicode=1
\pdfglyphtounicode{parenleftbig}{0028}
\pdfglyphtounicode{parenrightbig}{0029}
\pdfglyphtounicode{bracketleftbig}{005B}
\pdfglyphtounicode{bracketrightbig}{005D}
\pdfglyphtounicode{ceilingleftbig}{2308}
\pdfglyphtounicode{ceilingrightbig}{2309}
\pdfglyphtounicode{vextendsingle}{007C}
\pdfglyphtounicode{parenleftbigg}{0028}
\pdfglyphtounicode{parenrightbigg}{0029}
\pdfglyphtounicode{circleplustext}{2A01}
\pdfglyphtounicode{circleplusdisplay}{2A01}
\pdfglyphtounicode{summationtext}{2211}
\pdfglyphtounicode{producttext}{220F}
\pdfglyphtounicode{logicalandtext}{22C0}
\pdfglyphtounicode{summationdisplay}{2211}
\pdfglyphtounicode{productdisplay}{220F}
\pdfglyphtounicode{uniondisplay}{22C3}
\pdfglyphtounicode{hatwide}{0302}
\pdfglyphtounicode{hatwider}{0302}
\pdfglyphtounicode{tildewide}{0303}
\pdfglyphtounicode{mapsto}{007C}
\pdfglyphtounicode{axisshort}{002D}
\pdfglyphtounicode{arrowaxisright}{2192}
\usepackage{amsmath,amssymb,amsthm}
\newcommand{\mathscr}{\mathcal}
\usepackage{mathtools,microtype,enumitem,booktabs,graphicx}
\usepackage{tikz}
\usetikzlibrary{arrows.meta,positioning}
\usepackage[hidelinks,unicode,hypertexnames=false,bookmarksdepth=2]{hyperref}
\hypersetup{pdftitle={Minimal models and Mori fibre spaces for generalized log canonical Q-pairs},
 pdfauthor={OpenAI},pdfsubject={Existence of minimal models and Mori fibre spaces with fixed rational nef data}}
\usepackage[capitalise,nameinlink,noabbrev]{cleveref}
\setlength{\textwidth}{6.45in}
\setlength{\oddsidemargin}{0.025in}
\setlength{\evensidemargin}{0.025in}
\setlength{\textheight}{9in}
\setlength{\topmargin}{-0.2in}
\setlength{\headheight}{12pt}
\setlength{\headsep}{20pt}
\setlength{\emergencystretch}{2em}
\numberwithin{equation}{section}
\newtheorem{theorem}{Theorem}[section]
\newtheorem{proposition}[theorem]{Proposition}
\newtheorem{lemma}[theorem]{Lemma}
\newtheorem{corollary}[theorem]{Corollary}
\newtheorem{claim}[theorem]{Claim}
\theoremstyle{definition}
\newtheorem{definition}[theorem]{Definition}
\newtheorem{notation}[theorem]{Notation}
\newtheorem{example}[theorem]{Example}
\theoremstyle{remark}
\newtheorem{remark}[theorem]{Remark}
\newcommand{\C}{\mathbb C}
\newcommand{\Q}{\mathbb Q}
\newcommand{\R}{\mathbb R}
\newcommand{\Z}{\mathbb Z}
\newcommand{\N}{\mathbb N}
\newcommand{\PP}{\mathbb P}
\newcommand{\cO}{\mathcal O}
\DeclareMathOperator{\Spec}{Spec}
\DeclareMathOperator{\Proj}{Proj}
\DeclareMathOperator{\Supp}{Supp}
\DeclareMathOperator{\Pic}{Pic}
\DeclareMathOperator{\Cl}{Cl}
\DeclareMathOperator{\Hom}{Hom}
\DeclareMathOperator{\ord}{ord}
\DeclareMathOperator{\phaseangle}{angle}
\DeclareMathOperator{\vol}{vol}
\DeclareMathOperator{\lct}{lct}
\DeclareMathOperator{\gr}{gr}
\DeclareMathOperator{\rk}{rk}
\DeclareMathOperator{\relint}{relint}
\DeclareMathOperator{\Frac}{Frac}
\DeclareMathOperator{\trdeg}{trdeg}
\DeclareMathOperator{\codim}{codim}
\DeclareMathOperator{\mult}{mult}
\DeclareMathOperator{\Exc}{Exc}
\DeclareMathOperator{\id}{id}
\DeclareMathOperator{\im}{im}
\DeclareMathOperator{\Tr}{Tr}
\setlist[enumerate]{label=\textup{(\roman*)},leftmargin=*}
\setlist[itemize]{leftmargin=*}
